% =============================================================================
% Paper 8: Towards Standardization in Quantum Programming -- A Survey of
%          Frameworks, Intermediate Representations, and Interoperability
% =============================================================================
\documentclass[journal]{IEEEtran}

% --- Packages ----------------------------------------------------------------
\usepackage{cite}
\usepackage{amsmath,amssymb,amsfonts}
\usepackage{graphicx}
\usepackage{textcomp}
\usepackage{xcolor}
\usepackage{listings}
\usepackage{booktabs}
\usepackage{multirow}
\usepackage{hyperref}
\usepackage{tikz}
\usepackage{subcaption}
\usepackage{balance}
\usepackage{enumitem}
\usepackage{tabularx}
\usepackage{longtable}
\usepackage{url}

\usetikzlibrary{shapes, arrows.meta, positioning, fit, backgrounds, calc}

% --- Listings Configuration --------------------------------------------------
\lstset{
  basicstyle=\ttfamily\footnotesize,
  keywordstyle=\color{blue}\bfseries,
  commentstyle=\color{gray}\itshape,
  stringstyle=\color{red},
  numbers=left,
  numberstyle=\tiny\color{gray},
  stepnumber=1,
  numbersep=5pt,
  frame=single,
  breaklines=true,
  captionpos=b,
  tabsize=2,
  showstringspaces=false,
}

% --- Title -------------------------------------------------------------------
\title{Towards Standardization in Quantum Programming: A Survey of Frameworks, Intermediate Representations, and Interoperability Tools}

\author{
  Umer~Farooq,
  Hussan~Waseem~Syed,
  Muhammad~Irtaza~Khan,
  Imran~Ashraf,~\IEEEmembership{Senior Member,~IEEE,}
  and~Muhammad~Nouman~Noor
  \thanks{U.~Farooq, H.~W.~Syed, and M.~I.~Khan are with the Department of Computer Science, FAST National University of Computer and Emerging Sciences, Islamabad, Pakistan (e-mail: \{i220891, i220893, i220911\}@nu.edu.pk).}
  \thanks{I.~Ashraf and M.~N.~Noor are with the Faculty of Computing, FAST National University of Computer and Emerging Sciences, Islamabad, Pakistan.}
  \thanks{This work was conducted under the Open Quantum Workbench initiative at FAST University.}
}

\begin{document}
\maketitle

% =============================================================================
% ABSTRACT
% =============================================================================
\begin{abstract}
The quantum computing software ecosystem has grown rapidly over the past decade, with major technology companies and research institutions developing proprietary frameworks, each with distinct programming models, gate libraries, and execution paradigms. This fragmentation---analogous to the early days of classical computing before standardized instruction sets and high-level languages emerged---creates significant barriers to code portability, reproducible research, collaborative development, and educational accessibility. In this survey, we provide a comprehensive analysis of the current quantum programming landscape, covering seven major frameworks (Qiskit, Cirq, PennyLane, PyQuil, Amazon Braket, Q\#, and Strawberry Fields), four intermediate representations (OpenQASM 2.0/3.0, Quil, QIR, and XACC IR), and five interoperability tools (t$|$ket$\rangle$, Mitiq, QCanvas, Orquestra, and Strangeworks). We systematically compare these systems across twelve dimensions including programming paradigm, gate library coverage, classical control flow support, hardware access, simulation capabilities, and community adoption. We analyze the role of OpenQASM~3.0 as a potential universal standard, examining its strengths, adoption challenges, and implementation gaps. Drawing on our experience developing QCanvas---a multi-framework compilation platform that uses OpenQASM~3.0 as a ``Rosetta Stone'' intermediate representation---we identify key challenges in cross-framework transpilation, including semantic differences in gate definitions, incompatible parameter conventions, and divergent measurement models. We conclude with a roadmap for achieving greater standardization, proposing concrete steps for the quantum computing community including standardized gate semantics, unified circuit metadata, and benchmark suites for interoperability validation.
\end{abstract}

\begin{IEEEkeywords}
quantum computing, programming frameworks, standardization, OpenQASM 3.0, interoperability, survey, quantum software engineering
\end{IEEEkeywords}

% =============================================================================
% I. INTRODUCTION
% =============================================================================
\section{Introduction}
\label{sec:introduction}

\IEEEPARstart{T}{he} field of quantum computing has undergone remarkable growth since the first small-scale quantum processors became available for cloud access in the mid-2010s. IBM's deployment of a 5-qubit processor on the cloud in 2016~\cite{ibmquantum}, followed by Google's quantum supremacy demonstration in 2019~\cite{arute2019quantum}, and continuous scaling of qubit counts across multiple platforms have accelerated the development of quantum software tools. Today, a researcher or developer wishing to program a quantum computer faces an abundance of choices: Qiskit~\cite{qiskit2024} from IBM, Cirq~\cite{cirq2024} from Google, PennyLane~\cite{pennylane2024} from Xanadu, PyQuil~\cite{smith2016practical} from Rigetti, Amazon Braket~\cite{braket2024} from AWS, Q\#~\cite{svore2018q} from Microsoft, and several others.

While this diversity reflects healthy competition and innovation, it also creates a \emph{fragmentation problem} that increasingly impedes progress:

\begin{itemize}
    \item \textbf{Code portability:} A quantum algorithm implemented in Qiskit cannot be directly executed using Cirq or PennyLane without manual rewriting.
    \item \textbf{Reproducibility:} Research papers that provide code in a single framework are difficult to reproduce by groups using different tools.
    \item \textbf{Educational barriers:} Students must learn multiple, incompatible APIs to gain broad quantum computing literacy.
    \item \textbf{Vendor lock-in:} Choosing a framework often implies commitment to a specific hardware vendor's ecosystem.
    \item \textbf{Duplication of effort:} The community maintains parallel implementations of the same quantum algorithms across multiple frameworks.
\end{itemize}

The classical computing world faced analogous challenges in its early decades, eventually converging on standard instruction sets (x86, ARM), intermediate representations (LLVM IR), and high-level languages (C, Python). The quantum computing community is beginning a similar standardization journey, with OpenQASM~3.0~\cite{cross2022openqasm} emerging as the most prominent candidate for a universal quantum intermediate representation.

This survey makes the following contributions:

\begin{enumerate}
    \item A \textbf{comprehensive taxonomy} of quantum programming frameworks, intermediate representations, and interoperability tools, systematically compared across twelve dimensions.
    \item A \textbf{detailed analysis of OpenQASM~3.0} as a standardization vehicle, including its feature coverage, adoption status, and implementation challenges.
    \item \textbf{Empirical insights from QCanvas}, a multi-framework compilation platform that we developed, identifying concrete cross-framework transpilation challenges.
    \item A \textbf{standardization roadmap} with actionable recommendations for the quantum computing community.
\end{enumerate}

The remainder of this paper is organized as follows. Section~\ref{sec:methodology} describes our survey methodology. Section~\ref{sec:frameworks} surveys quantum programming frameworks. Section~\ref{sec:irs} analyzes intermediate representations. Section~\ref{sec:tools} reviews interoperability tools. Section~\ref{sec:openqasm} provides a deep dive into OpenQASM~3.0. Section~\ref{sec:challenges} examines cross-framework challenges using QCanvas as a case study. Section~\ref{sec:roadmap} proposes a standardization roadmap, and Section~\ref{sec:conclusion} concludes.

% =============================================================================
% II. SURVEY METHODOLOGY
% =============================================================================
\section{Survey Methodology}
\label{sec:methodology}

\subsection{Scope and Selection Criteria}

We surveyed quantum programming tools that satisfy the following criteria:
\begin{enumerate}
    \item Actively maintained with releases or commits within the past 24 months (as of early 2026).
    \item Supports gate-based quantum circuit programming.
    \item Provides either a programming framework, an intermediate representation, or an interoperability tool.
    \item Publicly documented with available source code or comprehensive documentation.
\end{enumerate}

\subsection{Comparison Dimensions}

We evaluate each system across twelve dimensions:
\begin{enumerate}
    \item Programming paradigm (imperative, functional, declarative, visual)
    \item Primary language and API design
    \item Gate library coverage (single-qubit, multi-qubit, parameterized)
    \item Gate modifier support (control, inverse, power)
    \item Classical control flow (conditionals, loops, functions)
    \item Type system (classical types, quantum types)
    \item Measurement model (mid-circuit, terminal, conditional)
    \item Simulation capabilities (statevector, density matrix, noise)
    \item Hardware access (which QPU providers)
    \item QASM export/import support
    \item Community and ecosystem (GitHub stars, contributors, papers)
    \item Educational resources (tutorials, documentation quality)
\end{enumerate}

\subsection{Data Sources}

Information was gathered from official documentation, GitHub repositories, published papers, and hands-on experimentation with each framework.

% =============================================================================
% III. QUANTUM PROGRAMMING FRAMEWORKS
% =============================================================================
\section{Quantum Programming Frameworks}
\label{sec:frameworks}

\subsection{IBM Qiskit}

Qiskit~\cite{qiskit2024} is IBM's open-source quantum computing SDK, first released in 2017. It is the most widely adopted quantum framework by community size.

\textbf{Programming Model:} Imperative, object-oriented. Circuits are \texttt{QuantumCircuit} objects; gates are appended via method calls (\texttt{qc.h(0)}, \texttt{qc.cx(0,1)}).

\textbf{Architecture:} Qiskit comprises several components: Terra (core), Aer (simulators), Ignis (error mitigation, deprecated), and the Qiskit Runtime for cloud execution. Recent versions (1.0+) have consolidated into a single package.

\textbf{Gate Library:} Comprehensive, including standard gates (H, X, Y, Z, S, T, CNOT, CZ, SWAP, Toffoli), rotation gates (Rx, Ry, Rz, P, U), and custom gate definitions. Supports gate modifiers via explicit controlled and inverse gate variants (\texttt{ch}, \texttt{ccx}, \texttt{sdg}).

\textbf{Classical Control:} Supports \texttt{if\_test} for conditional operations based on classical registers. Qiskit 1.0 introduced improved classical control flow via the \texttt{classical\_function} module.

\textbf{QASM Support:} Native export to OpenQASM~2.0 via \texttt{qc.qasm()}. OpenQASM~3.0 export is available through the \texttt{qiskit-qasm3-import} package but coverage is incomplete.

\textbf{Hardware:} Direct access to IBM Quantum processors (Eagle 127-qubit, Heron 133-qubit, and upcoming systems).

\textbf{Community:} $>$15,000 GitHub stars, $>$600 contributors, extensive textbook~\cite{wootton2021teaching}.

\subsection{Google Cirq}

Cirq~\cite{cirq2024} is Google's quantum computing framework, designed for near-term quantum devices and noise-aware circuit design.

\textbf{Programming Model:} Functional-declarative. Qubits are explicit objects (\texttt{cirq.LineQubit}, \texttt{cirq.GridQubit}); circuits are lists of \texttt{Moment} objects containing operations.

\textbf{Gate Library:} Standard gates with uppercase naming (H, X, Y, Z, S, T, CNOT, CZ, SWAP). Supports gate powers (\texttt{cirq.X(q)**0.5}), controlled gates (\texttt{.controlled\_by()}), and custom unitary gates.

\textbf{Classical Control:} Limited built-in classical control flow. Classical conditions typically handled in Python wrapper code rather than in the circuit representation.

\textbf{QASM Support:} Export to OpenQASM~2.0 via \texttt{cirq.QasmOutput}. No native OpenQASM~3.0 support.

\textbf{Hardware:} Access to Google's Sycamore and future processors via Google Cloud.

\textbf{Community:} $>$4,000 GitHub stars, focused on NISQ algorithms and device characterization.

\subsection{Xanadu PennyLane}

PennyLane~\cite{pennylane2024} is Xanadu's framework for differentiable quantum computing and quantum machine learning.

\textbf{Programming Model:} Decorator-based functional paradigm. Quantum circuits are Python functions decorated with \texttt{@qml.qnode}, enabling automatic differentiation.

\textbf{Gate Library:} Standard gates with full-name conventions (\texttt{Hadamard}, \texttt{PauliX}, \texttt{CNOT}). Extensive template library for variational circuits (StronglyEntanglingLayers, BasicEntanglerLayers). Supports controlled variants (CY, CH, CRX, CRY, CRZ, CSWAP, CCZ).

\textbf{Classical Control:} Supports mid-circuit measurements and classical conditioning in recent versions. Rich integration with NumPy, PyTorch, JAX, and TensorFlow for hybrid quantum-classical optimization.

\textbf{QASM Support:} Import from OpenQASM~2.0 via \texttt{qml.from\_qasm}. No native OpenQASM~3.0 export.

\textbf{Hardware:} Plugin-based access to IBM, AWS, Google, IonQ, Rigetti, and Xanadu's photonic hardware.

\textbf{Community:} $>$5,000 GitHub stars, strong presence in quantum machine learning research.

\subsection{Rigetti PyQuil}

PyQuil~\cite{smith2016practical} uses the Quil instruction language and provides access to Rigetti's superconducting quantum processors.

\textbf{Programming Model:} Imperative, instruction-based. Programs are sequences of Quil instructions.

\textbf{Gate Library:} Standard gates plus Rigetti-specific native gate sets optimized for their hardware.

\textbf{QASM Support:} Uses Quil as its native IR rather than OpenQASM. Third-party tools exist for QASM conversion.

\textbf{Hardware:} Rigetti Ankaa-class processors via Rigetti Quantum Cloud Services.

\subsection{Amazon Braket}

Amazon Braket~\cite{braket2024} is AWS's managed quantum computing service providing access to multiple hardware vendors through a unified API.

\textbf{Programming Model:} Imperative, similar to Qiskit. Circuits are \texttt{Circuit} objects with gate methods.

\textbf{Gate Library:} Standard gates following conventions similar to Qiskit/Cirq.

\textbf{Hardware:} Multi-vendor access: IonQ (trapped ion), Rigetti (superconducting), OQC (superconducting), QuEra (neutral atom).

\textbf{QASM Support:} Supports OpenQASM~3.0 as an input format (one of the few to do so), though generation is limited.

\subsection{Microsoft Q\#}

Q\#~\cite{svore2018q} is Microsoft's domain-specific language for quantum programming, integrated with the Azure Quantum platform.

\textbf{Programming Model:} Standalone compiled language (not embedded in Python) with strong typing, pattern matching, and first-class quantum operations.

\textbf{Gate Library:} Standard gates as language-level operations. Custom operations defined as Q\# functions.

\textbf{QASM Support:} Q\# uses QIR (Quantum Intermediate Representation) based on LLVM IR rather than OpenQASM.

\textbf{Hardware:} Azure Quantum: IonQ, Quantinuum, Rigetti, Pasqal.

\subsection{Xanadu Strawberry Fields}

Strawberry Fields~\cite{killoran2019strawberry} is Xanadu's framework for photonic quantum computing.

\textbf{Programming Model:} Photonic-specific with continuous-variable operations (displacement, squeezing, beam splitters).

\textbf{Note:} As a continuous-variable framework, Strawberry Fields represents a fundamentally different paradigm from gate-based frameworks and is included for completeness.

\subsection{Framework Comparison Summary}

Table~\ref{tab:frameworkcomparison} provides a systematic comparison across key dimensions:

\begin{table*}[t]
\centering
\caption{Systematic Comparison of Quantum Programming Frameworks}
\label{tab:frameworkcomparison}
\begin{tabular}{@{}lccccccc@{}}
\toprule
\textbf{Dimension} & \textbf{Qiskit} & \textbf{Cirq} & \textbf{PennyLane} & \textbf{PyQuil} & \textbf{Braket} & \textbf{Q\#} & \textbf{S.F.} \\
\midrule
Language       & Python & Python & Python & Python & Python & Q\# & Python \\
Paradigm       & OOP & Func. & Decorator & Instr. & OOP & Compiled & Func. \\
Gate naming    & lower & UPPER & FullName & UPPER & lower & PascalCase & N/A \\
Param.\ convention & (angle, qubit) & gate(qubit)\textasciicircum{} & (wires=) & gate(angle) qubit & (qubit, angle) & op(qubit) & N/A \\
ctrl modifier  & Explicit & .controlled\_by() & Explicit & CONTROLLED & Explicit & Controlled & N/A \\
inv modifier   & sdg/tdg & **(-1) & adjoint() & DAGGER & Adjoint & Adjoint & N/A \\
QASM 2.0       & Export & Export & Import & --- & --- & --- & --- \\
QASM 3.0       & Partial & --- & --- & --- & Import & --- & --- \\
Native IR      & QASM & Cirq IR & PL IR & Quil & Braket IR & QIR & SF IR \\
Mid-circuit meas.  & \checkmark & \checkmark & \checkmark & \checkmark & \checkmark & \checkmark & N/A \\
Noise models   & \checkmark & \checkmark & \checkmark & \checkmark & \checkmark & \checkmark & \checkmark \\
GitHub stars   & 15k+ & 4k+ & 5k+ & 1.5k+ & 300+ & 9k+ & 700+ \\
\bottomrule
\end{tabular}
\end{table*}

% =============================================================================
% IV. INTERMEDIATE REPRESENTATIONS
% =============================================================================
\section{Intermediate Representations}
\label{sec:irs}

Intermediate representations (IRs) serve as the bridge between high-level quantum programs and low-level hardware instructions. The choice of IR profoundly affects portability, optimization opportunities, and hardware compatibility.

\subsection{OpenQASM 2.0}

OpenQASM~2.0~\cite{cross2017open} was released by IBM in 2017 and became the de facto standard for quantum circuit interchange. It provides:
\begin{itemize}
    \item Basic gate definitions (\texttt{gate}, \texttt{opaque}).
    \item Qubit and classical bit register declarations.
    \item Measurement and reset operations.
    \item Conditional execution (\texttt{if}).
    \item Include mechanism for standard gate library.
\end{itemize}

\textbf{Limitations:} No classical control flow beyond simple conditionals, no subroutines, no rich type system, no gate modifiers, no scoping. These limitations motivated the development of OpenQASM~3.0.

\subsection{OpenQASM 3.0}

OpenQASM~3.0~\cite{cross2022openqasm} represents a major overhaul, transforming QASM from a circuit description format into a full programming language:

\begin{itemize}
    \item \textbf{Type system:} \texttt{qubit}, \texttt{bit}, \texttt{int[n]}, \texttt{uint[n]}, \texttt{float[n]}, \texttt{angle[n]}, \texttt{bool}, \texttt{complex[float[n]]}, arrays.
    \item \textbf{Control flow:} \texttt{if/else}, \texttt{for}, \texttt{while}, \texttt{break}, \texttt{continue}.
    \item \textbf{Subroutines:} \texttt{def} with parameters, local variables, return values.
    \item \textbf{Gate modifiers:} \texttt{ctrl @}, \texttt{inv @}, \texttt{negctrl @}, \texttt{pow(k) @}.
    \item \textbf{Timing:} \texttt{delay}, \texttt{duration}, \texttt{box} (for hardware scheduling).
    \item \textbf{Pulse-level:} Calibration blocks, \texttt{defcal}, frames, ports (OpenPulse integration).
    \item \textbf{I/O:} \texttt{input}/\texttt{output} directives for parameterized circuits.
\end{itemize}

We provide a detailed analysis of OpenQASM~3.0 in Section~\ref{sec:openqasm}.

\subsection{Quil}

Quil~\cite{smith2016practical} (Quantum Instruction Language) is Rigetti's native IR, designed as a practical instruction set architecture for quantum computing:
\begin{itemize}
    \item Instruction-based format (one operation per line).
    \item Support for parameterized gates and classical memory.
    \item \texttt{DECLARE} for classical memory allocation.
    \item \texttt{DEFGATE} for custom unitary gates via matrix specification.
    \item \texttt{PRAGMA} directives for compiler hints.
\end{itemize}

\textbf{Comparison with QASM:} Quil is more hardware-oriented, with explicit memory management and pragma-based compiler directives. It lacks the rich type system and control flow of OpenQASM~3.0 but provides better native integration with Rigetti hardware.

\subsection{QIR (Quantum Intermediate Representation)}

QIR~\cite{qir2021} is Microsoft's IR based on LLVM IR, designed for the Azure Quantum ecosystem:
\begin{itemize}
    \item Built on LLVM's SSA (Static Single Assignment) form.
    \item Quantum operations represented as function calls into a quantum runtime.
    \item Full classical computation support inherited from LLVM.
    \item Target-agnostic by design.
\end{itemize}

\textbf{Comparison with QASM:} QIR provides superior classical computation support due to its LLVM foundation but has lower human readability and narrower community adoption compared to OpenQASM.

\subsection{XACC IR}

XACC (eXtreme-scale ACCelerator)~\cite{mccaskey2020xacc} provides a hardware-agnostic IR and plugin framework:
\begin{itemize}
    \item Graph-based intermediate representation.
    \item Plugin system for adding new hardware backends.
    \item Supports both gate-model and annealing-model quantum computers.
\end{itemize}

\subsection{IR Comparison Summary}

\begin{table}[t]
\centering
\caption{Comparison of Quantum Intermediate Representations}
\label{tab:ircomparison}
\begin{tabular}{@{}lccccc@{}}
\toprule
\textbf{Feature} & \textbf{QASM2} & \textbf{QASM3} & \textbf{Quil} & \textbf{QIR} & \textbf{XACC} \\
\midrule
Rich types    & --- & \checkmark & Partial & \checkmark & --- \\
Control flow  & if only & Full & Limited & Full & --- \\
Subroutines   & --- & \checkmark & --- & \checkmark & --- \\
Gate mods.    & --- & \checkmark & --- & --- & --- \\
Pulse-level   & --- & \checkmark & --- & --- & --- \\
Human-readable & \checkmark & \checkmark & \checkmark & --- & --- \\
LLVM-based    & --- & --- & --- & \checkmark & --- \\
Adoption      & High & Growing & Medium & Medium & Low \\
\bottomrule
\end{tabular}
\end{table}

% =============================================================================
% V. INTEROPERABILITY TOOLS
% =============================================================================
\section{Interoperability Tools}
\label{sec:tools}

Several tools have been developed to bridge the gaps between quantum frameworks:

\subsection{Quantinuum t$|$ket$\rangle$}

t$|$ket$\rangle$~\cite{sivarajah2020tket} is a quantum compiler and circuit optimization toolkit from Quantinuum (formerly Cambridge Quantum):
\begin{itemize}
    \item Supports circuit input from Qiskit, Cirq, PennyLane, PyQuil, and Q\#.
    \item Powerful optimization passes (gate synthesis, routing, noise-aware compilation).
    \item Outputs to multiple backend formats.
    \item Does \emph{not} accept source code---requires programmatic circuit objects.
\end{itemize}

\subsection{Mitiq}

Mitiq~\cite{larose2022mitiq} focuses on quantum error mitigation:
\begin{itemize}
    \item Framework-agnostic error mitigation techniques (ZNE, PEC, CDR).
    \item Interfaces with Qiskit, Cirq, PennyLane, and Braket circuits.
    \item Operates at the circuit level rather than source code level.
\end{itemize}

\subsection{QCanvas (Our System)}

QCanvas is our multi-framework compilation platform, offering unique capabilities:
\begin{itemize}
    \item \textbf{Source code input:} Accepts framework-specific Python source code (not just circuit objects).
    \item \textbf{AST-based parsing:} Secure transpilation without code execution.
    \item \textbf{OpenQASM~3.0 output:} Standards-compliant output with Iteration I and II feature support.
    \item \textbf{Web-based IDE:} Browser-accessible development environment.
    \item \textbf{Multi-backend simulation:} Execute via Cirq, Qiskit, or PennyLane simulators.
\end{itemize}

We use QCanvas as a case study for cross-framework transpilation challenges in Section~\ref{sec:challenges}.

\subsection{Other Tools}

\textbf{Orquestra} (Zapata Computing) provides a workflow platform for quantum computing with multi-framework support but focuses on workflow orchestration rather than code-level interoperability.

\textbf{Strangeworks} provides a cloud platform aggregating multiple quantum hardware providers but does not address source code transpilation.

\subsection{Interoperability Tool Comparison}

\begin{table}[t]
\centering
\caption{Comparison of Quantum Interoperability Tools}
\label{tab:toolcomparison}
\begin{tabular}{@{}lccccc@{}}
\toprule
\textbf{Feature} & \textbf{t$|$ket$\rangle$} & \textbf{Mitiq} & \textbf{QCanvas} & \textbf{Orq.} & \textbf{Strgw.} \\
\midrule
Source code input   & --- & --- & \checkmark & --- & --- \\
AST-based parsing   & --- & --- & \checkmark & --- & --- \\
Circuit optimization & \checkmark & --- & Partial & --- & --- \\
Error mitigation    & --- & \checkmark & --- & \checkmark & --- \\
QASM 3.0 output     & Partial & --- & \checkmark & --- & --- \\
Web IDE             & --- & --- & \checkmark & Partial & \checkmark \\
Open source         & \checkmark & \checkmark & \checkmark & Partial & --- \\
\bottomrule
\end{tabular}
\end{table}

% =============================================================================
% VI. OPENQASM 3.0: DEEP DIVE
% =============================================================================
\section{OpenQASM 3.0 as a Standardization Vehicle}
\label{sec:openqasm}

OpenQASM~3.0 is the most promising candidate for a universal quantum programming standard. In this section, we provide a detailed analysis of its capabilities, adoption status, and challenges.

\subsection{Language Feature Analysis}

The OpenQASM~3.0 specification is extensive, containing features that span from basic quantum instructions to pulse-level hardware control. Based on our implementation experience with QCanvas, we categorize features by practical implementation difficulty:

\textbf{Core features (straightforward to implement):}
\begin{itemize}
    \item Version string, includes, comments.
    \item Qubit and bit declarations.
    \item Standard gate application with parameters.
    \item Measurement, reset, barrier.
    \item Basic arithmetic and comparison operations.
    \item If/else statements and for loops.
\end{itemize}

\textbf{Intermediate features (moderate difficulty):}
\begin{itemize}
    \item Gate modifiers (\texttt{ctrl@}, \texttt{inv@}).
    \item Complete type system (\texttt{int}, \texttt{uint}, \texttt{float}, \texttt{angle}).
    \item Arrays, aliasing, index sets.
    \item Input/output directives.
    \item Mathematical functions and constants.
\end{itemize}

\textbf{Advanced features (significant implementation effort):}
\begin{itemize}
    \item \texttt{negctrl@} and \texttt{pow(k)@} modifiers.
    \item Subroutines with parameter passing and scoping.
    \item While loops with break/continue.
    \item Complex type support.
    \item Bitwise and shift operators.
\end{itemize}

\textbf{Specialized features (often deferred or excluded):}
\begin{itemize}
    \item Timing constructs (\texttt{delay}, \texttt{duration}, \texttt{box}).
    \item Pulse-level control (\texttt{defcal}, frames, ports).
    \item Extern functions.
    \item Physical qubits (\texttt{\$0}, \texttt{\$1}).
    \item Hardware-specific extensions.
\end{itemize}

\subsection{Adoption Status}

Despite being published in 2022, OpenQASM~3.0 adoption remains uneven:

\begin{itemize}
    \item \textbf{Qiskit:} Partial support through \texttt{qiskit-qasm3-import}; active development.
    \item \textbf{Cirq:} No native OpenQASM~3.0 support (only QASM~2.0 export).
    \item \textbf{PennyLane:} QASM~2.0 import only; no 3.0 support.
    \item \textbf{Amazon Braket:} Supports OpenQASM~3.0 as input format (notable early adopter).
    \item \textbf{t$|$ket$\rangle$:} Partial QASM~3.0 support in recent versions.
    \item \textbf{QCanvas:} Comprehensive QASM~3.0 generation (Iteration I and II features).
\end{itemize}

This adoption gap is one of the key motivations for tools like QCanvas: even as the standard exists, practical tooling for generating compliant QASM~3.0 from existing framework code remains scarce.

\subsection{Implementation Challenges}

From our QCanvas implementation experience, we identify several practical challenges in QASM~3.0 code generation:

\begin{enumerate}
    \item \textbf{Measurement format ambiguity:} The specification allows both \texttt{measure q -> c} (QASM~2.0 style) and \texttt{c = measure q} (QASM~3.0 style). Ensuring consistent output requires explicit formatting choices.

    \item \textbf{For loop variable types:} The specification uses \texttt{uint} for loop variables, but some parsers expect \texttt{int}. We found that \texttt{int} provides better compatibility.

    \item \textbf{Gate modifier composition:} The semantics of composed modifiers (e.g., \texttt{ctrl @ inv @ x}) require careful ordering and qubit allocation.

    \item \textbf{Classical bit counting:} Determining the correct number of classical bits requires analyzing all measurement operations, not just register declarations.

    \item \textbf{Standard library dependency:} The \texttt{stdgates.inc} file is not standardized across implementations, leading to inconsistencies in available gate names.
\end{enumerate}

% =============================================================================
% VII. CROSS-FRAMEWORK CHALLENGES: QCANVAS CASE STUDY
% =============================================================================
\section{Cross-Framework Transpilation Challenges}
\label{sec:challenges}

Our development of QCanvas---which converts Qiskit, Cirq, and PennyLane source code to OpenQASM~3.0---revealed several categories of cross-framework incompatibility:

\subsection{Naming Convention Divergence}

The same quantum gate has different names across frameworks:
\begin{itemize}
    \item Hadamard: \texttt{h} (Qiskit), \texttt{H} (Cirq), \texttt{Hadamard} (PennyLane)
    \item CNOT: \texttt{cx}/\texttt{cnot} (Qiskit), \texttt{CNOT}/\texttt{CX} (Cirq), \texttt{CNOT} (PennyLane)
    \item Toffoli: \texttt{ccx} (Qiskit), \texttt{CCX}/\texttt{TOFFOLI} (Cirq), \texttt{Toffoli} (PennyLane)
    \item Phase: \texttt{p} (Qiskit), \texttt{ZPowGate} (Cirq), \texttt{PhaseShift} (PennyLane)
\end{itemize}

While conceptually straightforward to map, the sheer number of gate variants (30+ per framework) and aliasing (Qiskit's \texttt{cx} = \texttt{cnot}) makes comprehensive mapping non-trivial.

\subsection{Parameter Convention Incompatibility}

Frameworks disagree on parameter ordering and semantics:
\begin{itemize}
    \item \textbf{Qiskit:} \texttt{qc.rx(angle, qubit)} --- angle first, then qubit.
    \item \textbf{Cirq:} \texttt{cirq.rx(rads=angle)(qubit)} --- keyword argument on gate, applied to qubit.
    \item \textbf{PennyLane:} \texttt{qml.RX(angle, wires=qubit)} --- angle first, qubit as keyword.
    \item \textbf{Cirq power convention:} \texttt{cirq.X(q)**0.5} means $\sqrt{X}$, using exponentiation rather than explicit gate names.
\end{itemize}

\subsection{Circuit Construction Paradigm Mismatch}

The three frameworks use fundamentally different patterns:

\begin{itemize}
    \item \textbf{Qiskit (method-based):} \texttt{qc.h(0)} --- gate as method on circuit object.
    \item \textbf{Cirq (function-based):} \texttt{cirq.H(q0)} --- gate as function applied to qubit object.
    \item \textbf{PennyLane (declarative):} \texttt{qml.Hadamard(wires=0)} --- gate as module-level call inside decorated function.
\end{itemize}

These differences require entirely different AST visitor implementations, as described in our transpilation paper.

\subsection{Qubit Addressing Differences}

\begin{itemize}
    \item \textbf{Qiskit:} Integer indices (\texttt{qc.h(0)}).
    \item \textbf{Cirq:} Qubit objects (\texttt{cirq.H(q0)}), which can be \texttt{LineQubit}, \texttt{GridQubit}, or \texttt{NamedQubit}.
    \item \textbf{PennyLane:} Wire indices via keyword (\texttt{wires=0}) or wire lists (\texttt{wires=[0,1]}).
\end{itemize}

Mapping Cirq's grid qubits (e.g., \texttt{cirq.GridQubit(3, 5)}) to linear QASM indices requires a non-trivial qubit mapping strategy.

\subsection{Measurement Model Differences}

\begin{itemize}
    \item \textbf{Qiskit:} \texttt{qc.measure(qubit, clbit)} --- explicit qubit-to-clbit mapping.
    \item \textbf{Cirq:} \texttt{cirq.measure(q0, q1, key='result')} --- key-based with implicit classical register.
    \item \textbf{PennyLane:} \texttt{qml.expval(qml.PauliZ(0))} or \texttt{qml.counts()} --- return-based measurements with observable specification.
\end{itemize}

PennyLane's observable-based measurement model is particularly challenging to map to QASM's computational-basis measurement model.

\subsection{Inverse and Modifier Semantics}

Frameworks handle gate inversion differently:
\begin{itemize}
    \item Qiskit uses explicit inverse gates: \texttt{sdg}, \texttt{tdg}.
    \item Cirq uses exponentiation: \texttt{cirq.S(q)**-1}.
    \item PennyLane uses \texttt{qml.adjoint()}: \texttt{qml.adjoint(qml.S)(wires=0)}.
\end{itemize}

All must map to OpenQASM~3.0's \texttt{inv @ s q[0];} notation.

% =============================================================================
% VIII. STANDARDIZATION ROADMAP
% =============================================================================
\section{Towards a Standardization Roadmap}
\label{sec:roadmap}

Based on our survey analysis and practical experience with QCanvas, we propose a roadmap for achieving greater standardization in quantum programming:

\subsection{Short-Term (1--2 Years)}

\begin{enumerate}
    \item \textbf{Standardized Gate Semantics:} The community should agree on a canonical set of gate names, parameter conventions, and decomposition rules, published as a formal appendix to the OpenQASM~3.0 specification. This would eliminate the ambiguity in gate mapping tables.

    \item \textbf{QASM~3.0 Compliance Test Suite:} A standardized test suite for verifying OpenQASM~3.0 parser and generator compliance, similar to LLVM's test suite. QCanvas's 105-test suite could serve as a starting point.

    \item \textbf{Framework Export APIs:} Major frameworks should provide native, well-tested OpenQASM~3.0 export functionality. Currently, only Amazon Braket provides notable QASM~3.0 support.

    \item \textbf{Standard Gate Library File:} The \texttt{stdgates.inc} file should be formally standardized and versioned, ensuring all implementations use identical gate definitions.
\end{enumerate}

\subsection{Medium-Term (2--4 Years)}

\begin{enumerate}
    \item \textbf{Unified Circuit Metadata Standard:} A standard format for circuit metadata (qubit count, gate counts, depth, connectivity) that all frameworks can produce, enabling fair cross-framework benchmarking.

    \item \textbf{Cross-Framework Benchmark Suite:} A curated set of benchmark quantum circuits (similar to SPEC benchmarks for classical computing) with reference implementations in all major frameworks and verified QASM~3.0 equivalents.

    \item \textbf{Interoperability Certification:} A formal process for certifying that tools correctly convert between frameworks, using the benchmark suite for validation.

    \item \textbf{Pulse-Level Standardization:} Extending OpenQASM~3.0's pulse-level features (currently underspecified in practice) with concrete implementations and vendor-neutral pulse descriptions.
\end{enumerate}

\subsection{Long-Term (4+ Years)}

\begin{enumerate}
    \item \textbf{Universal Quantum Compiler Infrastructure:} An LLVM-like compiler infrastructure for quantum computing, with standardized IRs, optimization passes, and backend targets. QIR and OpenQASM~3.0 could converge into such an infrastructure.

    \item \textbf{Hardware-Agnostic Deployment:} A ``compile once, run anywhere'' model where quantum programs in OpenQASM~3.0 can be automatically compiled and deployed to any supported hardware platform.

    \item \textbf{Quantum Software Engineering Standards:} Formal software engineering standards for quantum programs, including testing methodologies, documentation requirements, and code quality metrics.
\end{enumerate}

\subsection{Lessons from Classical Computing}

The classical computing ecosystem offers instructive parallels:

\begin{itemize}
    \item \textbf{C language standardization (1970s--1990s):} Multiple incompatible C dialects converged through ANSI/ISO standardization. A similar process is needed for quantum assembly languages.
    \item \textbf{LLVM (2000s):} A universal compiler infrastructure enabled multiple frontend languages and backend targets. The quantum ecosystem needs its LLVM equivalent.
    \item \textbf{OpenGL/Vulkan (1990s--2010s):} Standardized graphics APIs eliminated vendor lock-in while allowing hardware-specific optimization. Quantum hardware APIs could follow this model.
    \item \textbf{Docker/Kubernetes (2010s):} Containerization standardized deployment. Quantum workload management could benefit from similar standardization.
\end{itemize}

% =============================================================================
% IX. DISCUSSION
% =============================================================================
\section{Discussion}
\label{sec:discussion}

\subsection{The Case for OpenQASM 3.0}

Our analysis suggests that OpenQASM~3.0 is well-positioned to serve as the primary standardization vehicle for quantum programming, for several reasons:

\begin{enumerate}
    \item \textbf{IBM backing and community support:} As the creator of both OpenQASM and the most widely used quantum framework (Qiskit), IBM has significant influence in driving adoption.
    \item \textbf{Language richness:} OpenQASM~3.0's feature set (types, control flow, subroutines, modifiers, pulse-level) covers the needs of most quantum programming scenarios.
    \item \textbf{Human readability:} Unlike QIR (LLVM-based, machine-oriented), OpenQASM~3.0 is readable by humans, making it suitable for education and debugging.
    \item \textbf{Existing adoption:} Many tools already support QASM~2.0; upgrading to 3.0 is an incremental rather than revolutionary change.
\end{enumerate}

\subsection{Barriers to Standardization}

Several factors slow standardization progress:

\begin{itemize}
    \item \textbf{Competitive dynamics:} Framework developers have commercial incentives to maintain proprietary ecosystems.
    \item \textbf{Rapid evolution:} The field is evolving so quickly that standards risk becoming obsolete.
    \item \textbf{Hardware diversity:} Different qubit technologies (superconducting, trapped ion, photonic, neutral atom) have fundamentally different native operation sets.
    \item \textbf{Specification completeness:} Parts of the OpenQASM~3.0 specification (particularly pulse-level features) remain underspecified.
\end{itemize}

\subsection{Role of Transpilation Tools}

Tools like QCanvas, t$|$ket$\rangle$, and Mitiq play a crucial bridging role during the transition period:
\begin{itemize}
    \item They demonstrate that cross-framework interoperability is \emph{feasible}.
    \item They identify practical standardization gaps through implementation experience.
    \item They provide immediate value to users while standards mature.
    \item They serve as reference implementations for how standards should be interpreted.
\end{itemize}

\subsection{Threats to Validity}

This survey has several limitations:
\begin{itemize}
    \item The quantum computing field evolves rapidly; some information may become outdated.
    \item Our practical experience is primarily with Qiskit, Cirq, and PennyLane; coverage of PyQuil, Braket, and Q\# relies more heavily on documentation analysis.
    \item Community adoption metrics (GitHub stars) are imperfect proxies for actual usage.
    \item The standardization roadmap reflects our assessment and may not align with industry priorities.
\end{itemize}

% =============================================================================
% X. CONCLUSION
% =============================================================================
\section{Conclusion}
\label{sec:conclusion}

The quantum computing software ecosystem faces a critical standardization challenge. Our survey of seven frameworks, four intermediate representations, and five interoperability tools reveals a fragmented landscape where code portability, reproducibility, and educational accessibility are significantly impaired by incompatible APIs, naming conventions, parameter semantics, and measurement models.

OpenQASM~3.0 emerges as the most promising standardization vehicle, offering a rich language that can serve as a universal intermediate representation. However, adoption remains uneven, and practical implementation challenges---as documented through our QCanvas case study---highlight the gap between specification and deployment.

We propose a concrete standardization roadmap organized into short-term (compliance test suites, framework export APIs), medium-term (benchmark suites, interoperability certification), and long-term (universal compiler infrastructure, hardware-agnostic deployment) objectives. Drawing parallels from classical computing's standardization history, we argue that the quantum community must actively pursue standardization to avoid prolonged fragmentation as the field scales.

The path forward requires collaboration between framework developers, standards bodies, hardware vendors, and the research community. Tools like QCanvas demonstrate that cross-framework interoperability is achievable today and can serve as both a bridge during the transition period and a reference for how standards should be practically implemented.

% =============================================================================
% ACKNOWLEDGMENT
% =============================================================================
\section*{Acknowledgment}

This work was conducted under the \emph{Open Quantum Workbench} initiative at FAST National University of Computer and Emerging Sciences. The authors thank the QSim team (Aneeq Ahmed Malik, Abeer Noor, Abdullah Mehmood) for contributions to the simulation backend.

% =============================================================================
% REFERENCES
% =============================================================================
\bibliographystyle{IEEEtran}

\begin{thebibliography}{30}

\bibitem{qiskit2024}
Qiskit contributors, ``Qiskit: An open-source framework for quantum computing,'' 2024. [Online]. Available: \url{https://qiskit.org/}

\bibitem{cirq2024}
Cirq Developers, ``Cirq: A Python library for writing, manipulating, and optimizing quantum circuits,'' Google AI Quantum, 2024.

\bibitem{pennylane2024}
V.~Bergholm \emph{et al.}, ``PennyLane: Automatic differentiation of hybrid quantum-classical computations,'' \emph{arXiv:1811.04968}, 2024.

\bibitem{cross2022openqasm}
A.~W.~Cross \emph{et al.}, ``OpenQASM~3: A broader and deeper quantum assembly language,'' \emph{ACM Transactions on Quantum Computing}, vol.~3, no.~3, pp.~1--50, 2022.

\bibitem{cross2017open}
A.~W.~Cross, L.~S.~Bishop, J.~A.~Smolin, and J.~M.~Gambetta, ``Open quantum assembly language,'' \emph{arXiv:1707.03429}, 2017.

\bibitem{ibmquantum}
IBM, ``IBM Quantum,'' 2024. [Online]. Available: \url{https://quantum.ibm.com/}

\bibitem{arute2019quantum}
F.~Arute \emph{et al.}, ``Quantum supremacy using a programmable superconducting processor,'' \emph{Nature}, vol.~574, no.~7779, pp.~505--510, 2019.

\bibitem{smith2016practical}
R.~S.~Smith, M.~J.~Curtis, and W.~J.~Zeng, ``A practical quantum instruction set architecture,'' \emph{arXiv:1608.03355}, 2016.

\bibitem{braket2024}
Amazon Web Services, ``Amazon Braket,'' 2024. [Online]. Available: \url{https://aws.amazon.com/braket/}

\bibitem{svore2018q}
K.~M.~Svore \emph{et al.}, ``Q\#: Enabling scalable quantum computing and development with a high-level DSL,'' in \emph{Proc.\ RWDSL}, 2018, pp.~7:1--7:10.

\bibitem{killoran2019strawberry}
N.~Killoran \emph{et al.}, ``Strawberry Fields: A software platform for photonic quantum computing,'' \emph{Quantum}, vol.~3, p.~129, 2019.

\bibitem{sivarajah2020tket}
S.~Sivarajah \emph{et al.}, ``t$|$ket$\rangle$: A retargetable compiler for NISQ devices,'' \emph{Quantum Science and Technology}, vol.~6, no.~1, p.~014003, 2020.

\bibitem{larose2022mitiq}
R.~LaRose \emph{et al.}, ``Mitiq: A software package for error mitigation on noisy quantum computers,'' \emph{Quantum}, vol.~6, p.~774, 2022.

\bibitem{qir2021}
QIR Alliance, ``Quantum Intermediate Representation (QIR) specification,'' 2021. [Online]. Available: \url{https://qir-alliance.org/}

\bibitem{mccaskey2020xacc}
A.~J.~McCaskey \emph{et al.}, ``XACC: A system-level software infrastructure for heterogeneous quantum-classical computing,'' \emph{Quantum Science and Technology}, vol.~5, no.~2, p.~024002, 2020.

\bibitem{wootton2021teaching}
J.~R.~Wootton, ``Teaching quantum computing with an interactive textbook,'' in \emph{Proc.\ IEEE QCE}, 2021, pp.~385--391.

\bibitem{fingerhuth2018open}
M.~Fingerhuth, T.~Babej, and P.~Wittek, ``Open source software in quantum computing,'' \emph{PLoS ONE}, vol.~13, no.~12, p.~e0208561, 2018.

\bibitem{alexander2020qiskit}
T.~Alexander \emph{et al.}, ``Qiskit pulse: Programming quantum computers through the cloud with pulses,'' \emph{Quantum Science and Technology}, vol.~5, no.~4, p.~044006, 2020.

\bibitem{barenco1995elementary}
A.~Barenco \emph{et al.}, ``Elementary gates for quantum computation,'' \emph{Physical Review A}, vol.~52, no.~5, pp.~3457--3467, 1995.

\bibitem{nielsen2010quantum}
M.~A.~Nielsen and I.~L.~Chuang, \emph{Quantum Computation and Quantum Information}.\hskip 1em plus 0.5em minus 0.4em Cambridge University Press, 10th anniversary ed., 2010.

\bibitem{peruzzo2014variational}
A.~Peruzzo \emph{et al.}, ``A variational eigenvalue solver on a photonic quantum processor,'' \emph{Nature Communications}, vol.~5, p.~4213, 2014.

\bibitem{farhi2014quantum}
E.~Farhi, J.~Goldstone, and S.~Gutmann, ``A quantum approximate optimization algorithm,'' \emph{arXiv:1411.4028}, 2014.

\bibitem{preskill2018quantum}
J.~Preskill, ``Quantum computing in the NISQ era and beyond,'' \emph{Quantum}, vol.~2, p.~79, 2018.

\bibitem{huang2019statistical}
Y.~Huang and M.~Martonosi, ``Statistical assertions for validating patterns and finding bugs in quantum programs,'' in \emph{Proc.\ ISCA}, 2019, pp.~541--553.

\bibitem{schuld2021machine}
M.~Schuld and F.~Petruccione, \emph{Machine Learning with Quantum Computers}.\hskip 1em plus 0.5em minus 0.4em Springer, 2nd ed., 2021.

\bibitem{nguyen2020pulse}
T.~Nguyen and A.~McCaskey, ``Enabling pulse-level programming, compilation, and execution in XACC,'' \emph{arXiv:2003.11971}, 2020.

\bibitem{mckay2018qiskit}
D.~C.~McKay \emph{et al.}, ``Qiskit backend specifications for OpenQASM and OpenPulse experiments,'' \emph{arXiv:1809.03452}, 2018.

\bibitem{wille2022tools}
R.~Wille, S.~Hillmich, and L.~Burgholzer, ``Tools for quantum computing based on decision diagrams,'' \emph{ACM Transactions on Quantum Computing}, vol.~3, no.~3, 2022.

\bibitem{steane1996error}
A.~M.~Steane, ``Error correcting codes in quantum theory,'' \emph{Physical Review Letters}, vol.~77, no.~5, pp.~793--797, 1996.

\bibitem{fu2017experimental}
X.~Fu \emph{et al.}, ``An experimental microarchitecture for a superconducting quantum processor,'' in \emph{Proc.\ MICRO}, 2017, pp.~813--825.

\end{thebibliography}

\end{document}
